\subsection*{11.2 Weak Law of Large Numbers}

In the followings, we consider a sequence of independent random variables X1,X2,

X3,.defined on a common probability space .(\Omega,\mathcal{F},P) F o r n\geq 1 , denote the

sum of the first n random variables by

Sn \coloneqqX 1 +X 2 +\cdot\cdot\cdot+ Xn.

We recall the following basic property of variance for pairwise uncorrelated

random variables (see Definition 7.3).

\begin{thmbox}{11.4}
\end{thmbox}

IfX1,X 2,...,X n are pairwise uncorrelated, then

Va r(X1 + X2 +\cdot\cdot\cdot+ Xn) =

n\sum

i=1

Va r(Xi).

\textit{Proof} Without loss of generality, suppose the random variables X1,...,X n have

zero mean. (We can consider the centered random variable Yn = Xn - E[Xn]

otherwise.)

Using the definition of variance and the linearity of expectation, we have

Va r(X1 + X2 +\cdot\cdot\cdot+ Xn) = E

[

X2

1 + X2

2 +\cdot\cdot\cdot+ X2

n + 2

\sum

i<j

XiXj

]

=

n\sum

i=1

E[X2

i ].

The assumption of pairwise uncorrelatedness guarantees that the cross-terms

E[XiXj ] are zero for i /=j Therefore, we have

Va r(X1 + X2 +\cdot\cdot\cdot+ Xn) =

n\sum

i=1

E[X2

i ]=

n\sum

i=1

Va r(Xi).

\hfill $\square$

We will present two versions of weak law of large numbers. The first one

assumes that the random variables are pairwise uncorrelated with the same mean

and variance. Note that the random variables need not be identically distributed, and

they need not be independent.

\begin{thmbox}{11.5 (Weak Law of Large Numbers (L2 Version))}
\end{thmbox}

Let .(Xi)\infty

i=1 be a sequence of pairwise uncorrelated random variables with

common mean E[Xi]= \mu and variance Va r(Xi) = \sigma2 for all i. Then

Sn/n

P

-\to \mu.

\textit{Proof} The proof is an application of Chebyshev inequality (Theorem 11.2).

Fix any \epsilon> 0 and integer n. Using the fact that Sn/n is an average of pairwise

uncorrelated random variables with common mean \mu,we have

P

(\vert\vertSn

n -\mu

\vert\vert >\epsilon

)

=P

(\vert\vertSn -n\mu

\vert\vert >n \epsilon

)

\leq Va r(Sn)

n2\epsilon2 = 1

n2\epsilon2

n\sum

i=1

Va r(Xi).

The last equality follows from Theorem 11.4.

Using Theorem 11.4, we know that .

\sumn

i=1 Va r(Xi)= n\sigma 2. Therefore, we have

P

( \vert\vert\vertSn

n -\mu

\vert\vert\vert >\epsilon

)

\leq n\sigma2

n2\epsilon2 = \sigma2

n\epsilon2 .

Asn\to\infty , the right-hand side approaches zero, which implies thatSn/n converges

in probability to \mu. This completes the proof. \hfill $\square$

The assumption of finite second moment can be relaxed to finite mean, but at the

cost of strengthening the uncorrelated assumption to independence.

\begin{thmbox}{11.6 (Weak Law of Large Numbers (L1 Version))}
\end{thmbox}

Suppose .(Xi)\infty

i=1 is a sequence of iid. random variables with finite mean \mu.

ThenSn/n

P

-\to \mu .

Theorem 11.5, also known as Khinchin's weak law of large numbers, has a

detailed proof in [4, Theorem 2.2.14]In the following two sections, we will discuss

two applications of this theorem instead of going through the proof.
